
\section{Identity and Determination}
\label{sec:identity-and-determination}

\subsection{The Ontological Problem of Determination}

The relation between identity and determination is fundamental to
understanding how an existent is the existent that it is. If being
and identity are ontologically inseparable, as established in the
preceding section, then the determinate character of an existent
cannot be treated as an accidental addition to an otherwise
complete and identity-free existence.

Yet this proposition raises a further question: what is the
relationship between identity and determination? Are they
different expressions of the same ontological condition, or
does determination constitute a broader field within which
identity is established?

The present analysis adopts a distinction between the general
condition of being determinate and the particular determinations
that contribute to the constitution of an existent's identity.

An existent is not an indeterminate actuality to which all
determinations are subsequently attached. Its existence is
already identity-bearing. Nevertheless, this does not imply
that every determination associated with it is constitutive
of its identity.

The distinction is essential. Without it, identity would become
indistinguishable from the totality of an entity's properties,
relations, states, and characteristics. Such an identification
would make identity excessively dependent upon every variation
in the existent and would eliminate the possibility of
distinguishing constitutive conditions from non-constitutive ones.

The problem, therefore, is not simply whether identity requires
determination. It is how determination relates to identity,
which determinations participate in its constitution, and
whether different determinations possess different ontological
roles.

\subsection{Determination as an Ontological Condition}

Determination concerns the fact that an existent is not merely
an unspecified possibility. It is constituted in a determinate
manner.

This formulation does not imply that determination must always
take the form of an externally observable property or an
independently identifiable feature. Nor does it imply that
every determination must be accessible through a particular
mode of manifestation.

Determination is considered here as an ontological condition
rather than merely a descriptive operation.

An existent may possess determinations that are not directly
manifested, not externally distinguishable, or not included
in any particular description. Their lack of accessibility
does not, by itself, establish their ontological absence.

It is therefore necessary to distinguish between determination
and the description of determination. The former concerns
the ontological constitution of an existent; the latter concerns
a representation or characterization of some aspect of that
constitution.

This distinction also prevents the reduction of determination
to classification. Assigning a category to an existent does
not necessarily establish the ontological conditions through
which it is determinate.

Determination must consequently be examined in relation to
the existent itself, rather than being defined exclusively
through the operations by which it may be described.

\subsection{Identity and the Constitutive Role of Determination}

Identity cannot be separated from the determinate actuality
of an existent. However, the relationship between identity
and determination requires more precision than the general
claim that identity involves determination.

Not every determination is necessarily constitutive of identity.
An existent may undergo changes in state, acquire or lose
particular relations, or exhibit different characteristics
without thereby becoming an entirely different existent.

If every determination were automatically identity-constitutive,
any change in any property or relation would necessarily entail
a change of identity. Such a conclusion cannot be assumed.

The present framework therefore distinguishes between
determinations that participate in constituting an existent's
identity and determinations that characterize other aspects
of its actuality.

A constitutive determination is not simply a feature that
happens to be associated with an entity. It contributes to
the ontological constitution through which that entity is
the entity that it is.

This does not yet provide a complete criterion for identifying
all constitutive determinations. Establishing such a criterion
requires further investigation into the structure of identity,
the conditions of persistence, and the relation between
determination and change.

At this stage, the distinction establishes a necessary
constraint: identity cannot be reduced to the aggregate of
all determinations, and the constitutive significance of
any particular determination must be demonstrated rather
than presumed.

\subsection{Internal and Relational Determination}

Determinations may be considered in terms of their different
ontological relations to an existent.

One important distinction concerns internal and relational
determination.

Internal determination refers to a determination considered
in relation to the constitution of the existent itself.
It concerns aspects of how the existent is constituted,
rather than merely how it is positioned in relation to
something else.

Relational determination concerns determinations that arise
through, or are constituted within, relations involving
the existent and other entities, conditions, or domains.

This distinction should not be interpreted as establishing
two completely independent classes of reality. An existent
may possess an internal constitution that is inseparable
from certain relations, and relational determinations may
have consequences for its internal constitution.

Nor should internal determination automatically be identified
with identity and relational determination with non-identity.
Such an identification would prematurely establish that
identity is exclusively intrinsic and that relations are
ontologically external to it.

The possibility that certain relations may be constitutive
of an existent's identity must remain open to investigation.

The relevant question is not whether a determination is
internal or relational in isolation, but what ontological
role it performs in the constitution of the existent.

\subsection{Structural Determination}

Structural determination concerns the organization and
configuration of determinations within an existent.

An existent may involve multiple determinations that are
not merely present as an unordered collection. Their
organization, interdependence, and relations may contribute
to the way in which the existent is constituted.

Structure must therefore be distinguished from the mere
aggregate of properties.

A collection of determinations does not necessarily constitute
a coherent structure. Structural determination involves
relations among determinations and the organization through
which they participate in a more comprehensive configuration.

The significance of structure for identity cannot, however,
be settled by assuming that every structural feature is
identity-constitutive.

Some structural arrangements may be essential to the identity
of an existent, while others may admit variation without
necessarily altering that identity. Determining which
arrangements are constitutive requires an account of the
specific ontological organization under consideration.

It is also necessary to avoid identifying structure with
signature. A structural configuration may contribute to
the characterization of an existent, but a description
of that configuration does not necessarily exhaust the
existent's identity.

The distinction between structural determination and
descriptive characterization will therefore remain important
throughout the subsequent analysis.

\subsection{Material Determination}

Material determination concerns determinations associated
with the material constitution of an existent.

Material constitution may contribute to the identity of
an existent in ways that cannot be reduced to a list of
observable material properties. The ontological significance
of material determinations depends upon the kind of existent
and the role that such determinations perform in its constitution.

It would therefore be premature to claim either that material
determinations always constitute identity or that they are
merely external conditions of its manifestation.

An existent may undergo material variation while retaining
some form of identity. Conversely, particular material
determinations may be indispensable to the identity of
certain existents.

These possibilities indicate that the relation between
material determination and identity cannot be resolved
through a universal assumption that every material change
either preserves or destroys identity.

The relevant conditions must be investigated in relation
to the ontological constitution of the existent and the
specific kind of identity under consideration.

Material determination is thus recognized as an ontologically
significant domain of inquiry, without being elevated into
a universal and exhaustive definition of identity.

\subsection{Determination and Distinction}

Determination is closely related to distinction, but the
two concepts must not be conflated.

To be determinate is to possess a specific ontological
constitution. To be distinct is to stand in a condition
of difference or non-identity in relation to something
else.

Determination may provide conditions through which an
existent is distinguishable, but distinguishability does
not necessarily exhaust its identity.

Two existents may share many determinations without becoming
numerically identical. Likewise, a particular existent
may possess determinations that do not distinguish it
from every other existent.

Consequently, the relationship between identity and
determination cannot be reduced to the capacity of a
description to differentiate one entity from another.

An account of identity must address the ontological
constitution of the existent itself, rather than merely
the conditions under which it can be distinguished.

This distinction also prevents the assumption that
uniqueness of a descriptive configuration is sufficient
to establish identity. A distinctive configuration may
provide evidence of difference, but its ontological
significance must be separately established.

\subsection{Determination and Change}

The relationship between determination and change presents
a particularly important problem.

An existent may undergo changes in some of its determinations
without necessarily losing its identity. If determination
were equivalent to identity in an undifferentiated sense,
every alteration would imply a corresponding identity
replacement.

However, the fact that an existent can change without
necessarily ceasing to be the existent that it is does
not establish that every change is compatible with the
preservation of identity.

Some changes may concern non-constitutive determinations.
Others may alter the organization or conditions through
which an identity is constituted. Still others may raise
the question of whether the existent remains numerically
the same existent.

The distinction between these possibilities cannot be
resolved by treating all determinations as ontologically
equivalent.

The present analysis therefore establishes a preliminary
distinction between the alteration of a determination
and the alteration of identity.

The former does not automatically entail the latter.
Nevertheless, whether a particular alteration affects
identity depends upon the constitutive role of the
determination involved.

This position avoids two opposite reductions: the assumption
that every change destroys identity, and the assumption
that identity remains unaffected by every possible change.

A more comprehensive account requires a separate examination
of identity, invariance, persistence, and continuity.

\subsection{Determination and the Limits of Identity}

The relation between determination and identity also
establishes limits on what may legitimately be claimed
about identity.

First, identity cannot be defined as the totality of all
determinations associated with an existent. Such a definition
would fail to distinguish constitutive determinations
from other determinations.

Second, identity cannot be reduced to a set of properties
selected for descriptive convenience. The selection of
features does not, by itself, establish their ontological
priority.

Third, no determination should be classified as
identity-constitutive merely because it is stable, observable,
or useful for distinguishing an existent.

Fourth, the absence of a determination from a particular
description does not establish its ontological absence.

Fifth, the possibility of relational and material
determinations participating in identity must remain
open rather than being excluded through an exclusively
intrinsic or immaterial conception of identity.

These limits are methodological and ontological constraints.
They do not constitute a complete theory of constitutive
determination.

The task of subsequent analysis is to establish more
precisely how the constitution of identity can be
distinguished from the broader field of determinations
belonging to an existent.

\subsection{Conceptual Distinctions in the Identity--Determination Relation}

The distinctions established in this section can be
summarized through the following conceptual relations:


\begin{table}[htbp]
    \centering
    \caption{Identity and Ontological Determination}
    \label{tab:identity-and-determination}

    \renewcommand{\arraystretch}{1.4}
    \begin{tabular}{p{3.2cm} p{5.1cm} p{5.3cm}}
        \hline
        \textbf{Dimension} &
        \textbf{Ontological Characterization} &
        \textbf{Implication for Identity} \\
        \hline

        Being and determination &
        Being is not an indeterminate substrate that first exists and subsequently receives determination. To be is to be ontologically determinate in some respect. &
        Identity is not an attribute added to an independently constituted entity. Being and identity are inseparable in the actuality of an entity. \\

        Constitutive determination &
        Determination concerns the definite character through which an entity is what it is. It belongs to the ontological constitution of the entity rather than merely to its external description. &
        Identity is inseparable from the determinations constitutive of the entity, although it cannot be reduced to an inventory of separately listed determinations. \\

        Determinate character &
        An entity is not an entirely indeterminate possibility. Its actuality involves a definite ontological character, even when that character is not exhaustively understood. &
        Identity presupposes determinateness, but the complete articulation of an entity's identity need not be available or reducible to a finite description. \\

        Internal determination &
        Some determinations belong to the constitution of an entity itself rather than arising solely from its relations to other entities. &
        The identity of an entity cannot be explained exclusively through its external relations or contextual descriptions. \\

        Relational determination &
        An entity may acquire or exhibit determinate conditions through its relations with other entities and domains. &
        Relational determinations may affect the entity's condition and expression without necessarily replacing or exhausting its identity. \\

        Structural determination &
        An entity may be constituted or constrained by an organized configuration of determinations and their relations. &
        Structural organization can contribute to identity, but identity should not be equated with a structural description alone. \\

        Material determination &
        An entity may possess determinations associated with its material constitution and realization. &
        Material characteristics may be constitutive of an entity at a particular ontological domain, but identity cannot be universally reduced to material properties. \\

        Determination and distinction &
        Determination establishes definite differences between an entity and what it is not, without requiring that every difference be identity-defining. &
        Not every distinguishing feature is constitutive of identity. The ontological significance of a determination must be established rather than assumed. \\

        Determination and change &
        Determinations may persist, vary, emerge, or cease through ontological processes. &
        Changes in determination raise the question of identity preservation, but neither every change destroys identity nor every persistence guarantees it. \\

        Determination and identity preservation &
        The continuity of an entity may involve the preservation of some determinations and the transformation of others. &
        Identity preservation requires an account of which determinations are constitutive and under what conditions their transformation remains compatible with the entity's identity. \\

        Limits of determination &
        No isolated determination necessarily provides an exhaustive account of an entity's ontological constitution. &
        Identity cannot be reduced to a single property, feature, relation, structural description, or material condition without further argument. \\

        \hline
    \end{tabular}
\end{table}


The table distinguishes the general ontological condition
of determination from its different forms and from the
specific question of constitutive significance for identity.

It should not be read as an exhaustive classification
of every possible determination. Nor does the placement
of a determination within a particular category establish
in advance whether it is identity-constitutive in every
existent.

Its purpose is to prevent the conceptual collapse of
identity into determination as a whole, while preserving
the necessary ontological relation between them.

\subsection{Concluding Formulation}

Identity and determination are ontologically related,
but they are not interchangeable.

An existent is not identity-bearing independently of
its determinate actuality. Nevertheless, the determinations
associated with an existent do not all necessarily
participate in constituting its identity.

The distinction between general determination and
constitutive determination is therefore indispensable
to the ontology of viability.

Identity concerns the determinate character through
which an existent is the existent that it is. Determination
concerns the broader ontological constitution through
which the existent is determinate in particular ways.

The constitutive significance of any determination
must be established through analysis of its ontological
role, rather than inferred merely from its presence,
stability, observability, or descriptive usefulness.

This section consequently establishes the following
principle:

\begin{quote}
Every existent is ontologically determinate, but not
every determination of an existent is necessarily
constitutive of its identity.
\end{quote}

This principle provides the basis for examining the
relation between identity and unity without reducing
either concept to the other.
